problem:  
Let \( (M^n, g) \) be a complete noncompact Riemannian manifold with nonnegative Ricci curvature. Assume its sectional curvature satisfies the borderline decay estimate  
\[
0 \le K(x) \le \frac{C}{(1+r(x))^2},
\]
where \(r(x) = d(x, x_0)\) is the distance from a fixed basepoint \(x_0\) and \(C > 0\) is a constant.  
Suppose \(M\) supports a nonconstant harmonic function \(u\) of linear growth, i.e.,  
\[
|u(x)| \le A(1 + r(x))
\]
for some \(A > 0\).

**Question:**  
Does the existence of such a linear-growth harmonic function necessarily imply that \(M\) contains a line, and hence (by the Cheeger–Gromoll theorem) splits isometrically as
\[
M = N^{n-1} \times \mathbb{R},
\]
with \(u\) coinciding, up to affine transformation, with the projection onto the \(\mathbb{R}\)-factor?  
Equivalently, under the borderline decay \(K(x) = O(r(x)^{-2})\), is every complete Ricci \(\ge 0\) manifold admitting a nontrivial linearly growing harmonic function globally split?

Why is it a "good" problem:  
This problem probes one of the most delicate thresholds in global Riemannian geometry: whether analytic information (existence of linear-growth harmonic functions) can force geometric rigidity (a global product decomposition).  

It directly builds upon the classical Cheeger–Gromoll splitting theorem and the extensive work by Li–Tam, Colding–Minicozzi, and Cheeger–Colding, but ventures beyond them by considering only borderline curvature decay \(K = O(r^{-2})\), a regime where neither standard comparison theorems nor asymptotic cone analysis straightforwardly apply.  

The \(r^{-2}\) decay rate is critical: stronger decay (say \(O(r^{-2-\varepsilon})\)) gives almost Euclidean behavior, while slower decay allows wild geometry at infinity. The question therefore targets the fine transition between rigidity and flexibility.  

Solving this would not only clarify the analytic–geometric interplay between harmonic functions and curvature decay, but might also yield new structural insights into the global geometry of nonnegatively Ricci curved manifolds and refine the borderline cases of Cheeger–Colding theory.  

In its conceptual simplicity—“Does a linear-growth harmonic function force a line under borderline curvature decay?”—the problem distills a fundamental geometric analysis question into a form inviting deep exploration across analysis, topology, and geometry.